\section{研究问题与关键实验（Research Questions）((要加结果和具体的以B为示例的实验步骤与分析))}
\label{sec:rq}

本节围绕“终局正确”与“机制正确”的可区分性，提出三个相互补充的研究问题（RQ），并给出对应的关键实验设计。三个 RQ 分别从\textbf{信息表达}、\textbf{机制扰动}与\textbf{动态学习}三个维度检验 LLM 的机制理解：RQ1 通过提示词信息密度梯度识别“参数堆叠式正确”与“机制推断式正确”；RQ2 通过结构参数的敏感性分析检验模型是否遵循理论比较静态；RQ3 通过多轮虚拟博弈检验模型是否能够在机制反馈下形成可解释的策略改进与收敛。三组实验均以 Scenario~1（推断外部性，主案例）为核心展示，并在 Scenario~2/3 上提供补充验证。

\subsection{RQ1：信息表达如何影响机制一致性？（提示词信息密度梯度）}
\label{subsec:rq1}

\paragraph{研究动机.}
现有评测常将“终局指标接近基准”视为模型理解机制的证据，但在复杂机制环境中，一个模型可能依赖参数提示、语言启发式或提示词诱导而“碰巧做对”。因此，我们提出一个\textbf{提示词信息密度梯度}实验：在保持经济环境、参数实例与理论基准不变的前提下，仅改变机制描述的表述方式与信息密度，以观察模型行为是否出现与机制理解一致的“拐点式改善”。

\paragraph{实验设计.}
我们在 Scenario~1 中构造同一机制的六种提示词版本 $\{\mathrm{v0},\ldots,\mathrm{v5}\}$，从简到繁逐步增加对机制结构的显式表达：
\begin{itemize}
  \item \textbf{v0--v2（规则/参数层）}: 仅给出任务规则与必要参数解释，考察模型对制度语义、变量含义与可执行行动的理解；
  \item \textbf{v3（机制概念注入）}: 显式引入“推断外部性”概念及其因果链条（相关性$\rightarrow$可推断性$\rightarrow$隐私损失$\rightarrow$价格/参与）；
  \item \textbf{v4（结构化机制解释）}: 在 v3 基础上给出更完整的机制分解与决策权衡框架，强调外部性如何进入个体的边际收益；
  \item \textbf{v5（公式解释框架）}: 进一步提供关键公式/符号解释，考察模型处理形式化表述与长文本结构的能力。
\end{itemize}
为控制混杂因素，每个版本在同一参数实例集上评测，并保持行动空间、交互协议与输出解析规则完全一致。对每个模型与提示词版本，我们进行 $K$ 次重复运行以估计均值与方差。

\paragraph{观测量与判据（机制拐点）.}
我们关注两类现象：（i）\textbf{单调堆叠 vs 机制拐点}：若模型表现随信息密度单调改善，可能反映其依赖“更多细节更容易猜对”；相反，若在 v3 处出现显著跃迁并在 v4/v5 边际收益递减，则更符合“机制概念注入触发推断”的解释；（ii）\textbf{结构一致性优先于终局一致性}：真正的机制理解应首先体现在结构对象（如分享集合结构、关键边际权衡方向）趋于一致，而不仅是终局指标数值接近。上述判据为后续“伪理解”提供可检验的诊断证据。

\subsection{RQ2：模型是否遵循机制比较静态？（敏感性分析）}
\label{subsec:rq2}

\paragraph{研究动机.}
机制理解不仅应在单一设置下给出合理行动，更应在机制参数变化时呈现与理论一致的比较静态响应。若模型只是“记住某个答案形态”或受提示诱导，其策略往往对关键参数扰动不稳定，甚至出现方向性错误。因此，我们对每个模型家族的代表性版本进行系统的敏感性分析，用以检验其行为是否符合机制的方向性预测。

\paragraph{扰动维度与设置.}
以 Scenario~1 为核心，我们围绕推断外部性链条选择三类关键扰动：
\begin{itemize}
  \item \textbf{相关性强度 $\rho$}: 提高 $\rho$ 应增强可推断性与外部性，从而改变“他人分享”对个体决策的边际影响；
  \item \textbf{隐私偏好分布 $\{v_i\}$}: 改变异质性与整体水平，用以检验模型是否能在个体层面做出一致的边际权衡（高 $v_i$ 更倾向拒绝，低 $v_i$ 更倾向参与，且该排序对机制扰动保持稳定）。
\end{itemize}
我们采用“同机制、跨设置”的参数网格或分层抽样生成实例集合，并在每个设置下重复运行以控制采样波动。

\paragraph{机制一致性的敏感性判据.}
敏感性分析不仅比较终局值的变化，更强调\textbf{方向性与结构稳定性}：当 $\rho$ 上升或噪声下降时，模型是否表现出更强的外部性意识（例如对“他人分享导致自身隐私损失”给出一致的解释与行动响应）；当 $v_i$ 分布平移或拉伸时，模型输出的个体行动排序是否保持单调一致。若模型在不同扰动下出现非稳定跳变或方向性相反的响应，则支持“伪理解”假设。该实验在 Scenario~2/3 上可对应扰动市场强度参数、匿名化约束强度与隐私成本分布，以验证结论的跨机制稳健性。

\subsection{RQ3：模型能否在机制反馈下学习并收敛？（多轮虚拟博弈）}
\label{subsec:rq3}

\paragraph{研究动机.}
静态一次性决策无法区分“偶然正确”与“可持续的机制学习”。在真实数据市场与平台治理中，主体通常在重复交互中观察反馈并更新策略。因此，我们引入\textbf{多轮虚拟博弈}作为机制学习测试：为模型提供由机制结算得到的历史反馈（而非人工点评），观察其是否能基于反馈形成可解释的策略改进，并逐步收敛至理论基准附近。该实验也为“弱模型在机制反馈下是否可被修复”提供证据。

\paragraph{虚拟博弈协议.}
在给定参数实例下，我们让同一模型连续进行 $T$ 轮交互。每一轮结束后，环境向模型反馈一组结构化摘要（例如：上一轮分享率/参与率、平台利润/福利分解、关键偏离来源的统计描述），并要求模型在下一轮基于该反馈更新其策略。为避免“泄露标准答案”，反馈仅包含由机制产生的\textbf{可观测结果}与\textbf{聚合统计}，不直接给出理论最优行动集合。

\paragraph{收敛与学习判据.}
我们用三类信号刻画机制学习：（i）\textbf{改进性}：随轮次推进，终局偏差与结构偏差是否显著下降；（ii）\textbf{稳定性}：策略是否趋于稳定而非振荡，且在重复运行下方差收缩；（iii）\textbf{可解释的机制一致理由}：模型在解释中是否越来越多引用外部性链条与边际权衡，而不是依赖模板化措辞。特别地，我们将该实验聚焦于在 RQ1 中对 v4 表现较弱的模型子集，以检验“机制反馈是否能补足其机制理解”。

\subsection{本节小结}
\label{subsec:rq_summary}

三个 RQ 共同构成从静态到动态、从语言到机制的分层验证体系：RQ1 在固定机制下改变信息表达以识别“参数≠理解”；RQ2 通过机制参数扰动检验比较静态一致性；RQ3 通过重复交互检验机制学习与收敛。后续实验结果将以统一的诊断指标体系呈现，并在主案例（Scenario~1）中重点报告提示词梯度的机制拐点、敏感性曲线的方向性一致性，以及虚拟博弈下的学习轨迹与收敛特征。
